
\documentclass[11pt,a4paper]{article}
\usepackage[a4paper,margin=2.05cm]{geometry}
\usepackage[spanish,es-noquoting]{babel}
\usepackage[T1]{fontenc}
\usepackage[utf8]{inputenc}
\usepackage{lmodern,microtype}
\usepackage{amsmath,amssymb,amsthm,mathtools}
\usepackage{booktabs,longtable,array,enumitem}
\usepackage{xcolor,hyperref}
\pagecolor{white}\color{black}
\setlength{\parindent}{0pt}
\setlength{\parskip}{0.65em}
\newtheorem{teorema}{Teorema}
\newtheorem{observacion}[teorema]{Observación}
\newtheorem{conjetura}[teorema]{Conjetura}
\title{\bfseries N41 --- CKM como cruce de canales HMT\\
\large \(A\), contraángulo, torsión y comparación PDG/CODATA}
\author{HMT / auditoría técnica}
\date{}
\begin{document}
\maketitle

\begin{abstract}
N41 calcula la matriz CKM HMT desde el ángulo eléctrico \(A=1000\alpha\), el paso de contraángulo \(h=C^\ast/6\) y la torsión angular \(\Delta^\circ\). Se compara con la actualización PDG 2025 de la matriz CKM y con CODATA 2018/2022 para \(\alpha\). El resultado es una coincidencia muy fina a nivel de fórmula angular/torsional. El documento separa el hecho numérico de la derivación completa por rutas: N41 certifica el cierre de fórmula y comparación; la derivación del species-pack de los coeficientes queda como siguiente capa.
\end{abstract}

\section{Datos HMT}

\[
\alpha_{\rm HMT}=0.007297352569283800660659977666000486.
\]
\[
A=1000\alpha=7.297352569283800605^\circ.
\]
\[
C^\ast=2.123738338968645945^\circ,\qquad h=\frac{C^\ast}{6}=0.353956389828107676^\circ.
\]
\[
\Delta_4=1.111964226921400488e-04\ {\rm rad},
\qquad
\Delta^\circ=0.006371085717212360^\circ.
\]

\section{Ángulos HMT}

\[
\boxed{\theta_{12}=2A-5h+28\Delta^\circ}
\]
\[
\boxed{\theta_{23}=7h-\frac{25}{2}\Delta^\circ}
\]
\[
\boxed{\theta_{13}=A-20h-\frac23\Delta^\circ}
\]
\[
\boxed{\delta_{\rm CKM}=9A}.
\]

Numéricamente:
\[
\theta_{12}=13.003313589509^\circ,
\quad
\theta_{23}=2.398056157332^\circ,
\quad
\theta_{13}=0.213977382244^\circ,
\quad
\delta_{\rm CKM}=65.676173123554^\circ.
\]

\section{Matriz CKM}

\[
|V_{\rm CKM}^{\rm HMT}|=
\begin{pmatrix}
0.974350259 & 0.225005836 & 0.003734601\\
0.224873110 & 0.973489279 & 0.041841465\\
0.008582400 & 0.041121745 & 0.999117283
\end{pmatrix}.
\]
\[
\boxed{J_{\rm HMT}=3.118972333e-05}.
\]

\section{Comparación con PDG 2025}

La actualización PDG 2025 da la matriz global-fit:
\[
|V|_{\rm PDG}=
\begin{pmatrix}
0.97435 & 0.22501 & 0.003732\\
0.22487 & 0.97349 & 0.04183\\
0.00858 & 0.04111 & 0.999118
\end{pmatrix}
\]
con las incertidumbres indicadas en el documento PDG.

En esta auditoría:
\[
\boxed{\max |\text{pull}| \text{ de los 9 módulos}=0.028902\sigma.}
\]
Además:
\[
\boxed{J_{\rm HMT}\text{ tiene pull }-0.008564\sigma.}
\]
\[
\boxed{\delta_{\rm HMT}\text{ tiene pull }-0.028251\sigma.}
\]

\section{Comparación CODATA de \(\alpha\)}

\[
\alpha_{\rm HMT}^{-1}=137.035999083999996628.
\]
CODATA 2018:
\[
\alpha^{-1}=137.035999084(21).
\]
CODATA 2022:
\[
\alpha^{-1}=137.035999177(21).
\]
Pulls:
\[
\boxed{\text{2018: }0.000000\sigma,}
\qquad
\boxed{\text{2022: }-4.428572\sigma.}
\]

\section{Lectura}

N41 muestra que el bloque CKM HMT no es una matriz pegada a mano. Sus tres ángulos mezclan:
\[
A=\text{cierre eléctrico},
\qquad
h=\text{paso de contraángulo},
\qquad
\Delta^\circ=\text{torsión angular}.
\]
La fase global:
\[
\delta=9A
\]
cierra sobre la rueda. En lectura HMT:
\[
\boxed{\text{CKM es cruce angular/torsional de canales de sección.}}
\]

\section{Disciplina}

N41 no prueba por sí solo el species-pack de rutas que produce los coeficientes
\[
28,\quad -25/2,\quad -2/3.
\]
Ese es el siguiente cierre. N41 prueba el estatuto de fórmula y comparación:
\[
\boxed{\text{con } A,h,\Delta^\circ \text{ fijados, la matriz cae dentro de PDG con precisión muy alta.}}
\]
\end{document}
